\section{Discussion: when is recurrent depth test-time compute?}
\label{sec:discussion}

The same architecture behaves as a parameter-efficient model on text and as an extendable serial computer on the
chain task. We offer a reading of the difference. It is an interpretation of two settings, not
a tested law: in this reading, recurrent depth works as test-time compute when three conditions hold.

\paragraph{Headroom.} The model's errors must come from too few serial steps rather than from missing knowledge. On
the chain task the dense models know every table and fail only on depth; one more loop is one more step. On web text at
these sizes we found no sign of such errors: the per-token gain is a diffuse shift rather than a concentrated one, and
extra loops at inference never helped, at any learning rate (Sections~\ref{sec:where} and~\ref{sec:mechanism}). This is
the hardest of the three conditions to measure independently, since here it is inferred from where extra loops do and
do not help. The larger gains on math and code text hint at more headroom there, but at this size those sets are
dominated by which documents a run happened to see.

\paragraph{Extendable iteration.} The loop must be able to run longer without destroying its state. Which kind of
iteration a loop learns is decided by its training more than by the architecture. On text the same block becomes a
fixed-length program, a periodic cycle or a fixed point, depending on the loop schedule (Section~\ref{sec:mechanism}).
On the chain task the loop schedule decides whether the learned algorithm survives extra loops: models trained with as
many loops as steps lose their values when run longer, and a few random extra loops in training fix that; the learning
rate decides whether the task is learned at all (Sections~\ref{sec:probes}--\ref{sec:stopping}). A fixed point is
extendable but, on text, computes nothing more; a fixed-length program computes but cannot be extended.

\paragraph{A stopping signal.} Extra loops are only useful if the model can tell how many to run. On the chain task the
model's own confidence is such a signal for most regimes, and training a few loops past the last step makes the
regime that extends furthest stoppable as well (Section~\ref{sec:stopping}). We did not test stopping rules on text,
where extra loops did not help.

The chain task meets all three conditions once the model is trained for them; web text at these sizes showed no sign
of headroom, and its one extendable regime, the fixed point, gained nothing from running longer. Two follow-ups we did
not run would test this reading: a serial task whose input length does not grow with the number of steps, which would
also decide whether a written chain of thought extends when the prompt does not get longer; and training on text that
needs more serial steps, such as mathematics and code, compared at equal compute.

\paragraph{The value of a loop is not a constant of the architecture.} Even where the loop does not iterate, it buys
something on text, and what it buys depends on size and data: at 10--20M with up to 20 tokens per parameter a loop is
worth at most about its parameters, at 40M--160M with 20--40 tokens per parameter about 1.3--1.5 times them.
Effective-parameter laws whose loop gain depends on the loop count but not on model size or training tokens describe
the second regime well and the first not at all, and the single pre-registered $\varphi$ averaged over both. Studies at one scale will undersell or oversell the loop
depending on where they sit.

\paragraph{Practical guidance.} Within the range we measured:
\begin{itemize}
  \item Loop when parameters or data are the constraint, not compute or latency. The advantage at equal parameters
  grows with model size and tokens per parameter; at equal training compute a dense model was better in every
  measured comparison but one, a tie at the smallest size and the most data with the looped model's own learning rate; at deployment the loop pays its executed layers in decoding time and KV cache.
  \item Loop a middle block with full backpropagation, start with four loops on text, and judge at the target scale:
  at 40M--80M eight loops did about as well as four at about 1.75 times the training compute, and small-scale ablations undersell
  the loop.
  \item Tune the looped model's learning rate separately, especially at small scale and on tasks that need serial
  steps: at a learning rate that suits the dense model, most looped runs did not learn the chain task in the given
  steps.
  \item Do not expect test-time scaling on text from standard training. To see whether a model can use more loops,
  sweep the loop count at inference; to see whether it computes step by step even when it cannot, probe it.
  \item If a model must decide by itself how long to run, its confidence is a good stopping signal, and training a few
  loops past the last step makes the stop reliable. On our task, a written chain of thought was the cheaper choice
  within the training range and failed beyond it; that failure came from longer prompts, so it says little about tasks
  whose input length is fixed.
\end{itemize}
